% Fichier généré automatiquement par tools/generate_corrections.py.
% Source : CIN/CIN-03-Transmetteurs/35_Vario/35_Vario.tex
% Statut : complété (2/2 question(s)).
\begin{corrigebox}[Corrigé complété]
\CorrigeQuestion{1}
On choisit les inconnues et les conventions positives de la figure,
        on écrit les relations de conservation/fermeture adaptées, puis on résout le
        système obtenu. Le résultat symbolique doit être homogène et vérifié dans les cas
        limites (entrée nulle, rapport unitaire ou position de référence).
\CorrigeQuestion{2}
On cherche une relation entre $\omega_{\text{Mh}/0}$, $\omega_{\text{Ph}/0}$ et $\omega_{\text{Mm}/0}$ (avec Mm et 4 même classe d'équivalence). Pour cela, on va d'abord rechercher une relation entre $\omega(3/0)$, $\omega(4/0)$ et $\omega(1/0)$.

Bloquons le porte satellite 4, directement lié au moteur Mm. On est alors en présence d'un réducteur simple d'entrée  $\omega(1/4)$ et de sortie $\omega(3/4)$. On a donc : 
$\dfrac{\omega(3/4)}{\omega(1/4)} = -\dfrac{R_{12}}{R_{32}}$. 

En libérant le porte satellite, on a donc :
$ \dfrac{\omega(3/4)}{\omega(1/4)}
= \dfrac{\omega(3/0)-\omega(4/0)}{\omega(1/0)-\omega(4/0)}
= -\dfrac{R_{12}}{R_{32}}
\Leftrightarrow 
R_{32} \omega(3/0) +R_{12}\omega(1/0) = \omega(4/0)\left(R_{12}+R_{32}\right)$

On a donc, $R_{32} \omega(3/0) +R_{12}\omega(1/0) = \omega(\text{Mm}/0)\left(R_{12}+R_{32}\right)$.

Par ailleurs, $\dfrac{\omega(\text{Ph}/0)}{\omega(3/0)} = -\dfrac{R_{3P}}{R_P}$ et 
$\dfrac{\omega(1/0)}{\omega(\text{Mh}/0)} = -\dfrac{R_M}{R_{1M}}$.

On a donc, $ \dfrac{2y}{x} \omega(Mh/0)= -\omega(3/0)\dfrac{R_{3P}}{R_P} \Leftrightarrow  \omega(3/0) = -\dfrac{2y}{x} \dfrac{R_P}{R_{3P}}\omega(Mh/0)$.

En utilisant la relation du train épi :
On a donc, $-R_{32} \dfrac{2y}{x} \dfrac{R_P}{R_{3P}}\omega(Mh/0)  -R_{12} \dfrac{R_M}{R_{1M}} \omega(\text{Mh}/0) = \omega(\text{Mm}/0)\left(R_{12}+R_{32}\right) \Leftrightarrow \left(-R_{32} \dfrac{2y}{x} \dfrac{R_P}{R_{3P}}  -R_{12} \dfrac{R_M}{R_{1M}} \right)\omega(\text{Mh}/0) = \omega(\text{Mm}/0)\left(R_{12}+R_{32}\right)$.

$\dfrac{\omega(Mh/0)}{\omega(Mm/0)}=-\dfrac{R_{12}+R_{32}}{R_{32} \dfrac{2y}{x} \dfrac{R_P}{R_{3P}}  +R_{12} \dfrac{R_M}{R_{1M}}}$

$\dfrac{\omega(Mh/0)}{\omega(Mm/0)}=-\dfrac{\left( R_{12}+R_{32}\right)R_{1M} R_{3P}x }{R_{32} 2y R_PR_{1M} + R_{3P}xR_{12} R_M}$. 
On a donc, $A=\left( R_{12}+R_{32}\right)R_{1M} R_{3P}$, $B=R_{32} 2R_{1M}$ et $C= R_{3P}xR_{12} R_M$. 
\textbf{Attention, plusieurs solutions possibles, si on factorise le numérateur et le dénominateur par l'un ou l'autre des rayons.}
\end{corrigebox}
